% Part: first-order-logic
% Chapter: tableaux
% Section: rules-and-proofs

\documentclass[../../../include/open-logic-section]{subfiles}

\begin{document}

\iftag{FOL}
      {\olfileid{fol}{tab}{rul}}
      {\olfileid{pl}{tab}{rul}}

\olsection{నియమాలు మరియు \usetoken{P}{tableau}}

ఒక \tetoken{టాబ్లో}{!!^a{tableau}} అనేది ఏదో ఒక \tetoken{నిర్మాణంలో}{!!a{structure}} ఒక \tetoken{వాక్యం}{!!a{sentence}}
సత్యం లేదా అసత్యం కాగల సాధ్యమైన మార్గాల క్రమబద్ధ పరిశీలన. టాబ్లో
నిర్మాణ ఘటకాలు \tetoken{చిహ్నిత సూత్రాలు}{!!{signed formula}s}: ఒక సత్యమూల్య “సంకేతం”
($\True$ లేదా~$\False$)తో కూడిన \tetoken{వాక్యాలు}{!!{sentence}s}. ఈ చిహ్నిత
\tetoken{సూత్రాలను}{!!{formula}s} కిందివైపుకు పెరిగే వృక్షంలో అమరుస్తాం.

\begin{defn}
  ఒక \emph{\tetoken{చిహ్నిత సూత్రం}{!!{signed formula}}} అంటే ఒక సత్యమూల్యం, ఒక
  \tetoken{వాక్యంతో}{!!a{sentence}} కూడిన జత; అంటే, కింది రెండింటిలో ఒకటి:
  \[
  \sFmla{\True}{!A} \text{ or } \sFmla{\False}{!A}.
  \]
\end{defn}

సహజార్థంలో $\sFmla{\True}{!A}$ను “$!A$ సత్యం కావచ్చు” అనీ,
$\sFmla{\False}{!A}$ను “$!A$ అసత్యం కావచ్చు” అనీ (ఏదో ఒక
\tetoken{నిర్మాణంలో}{!!{structure}}) చదవవచ్చు.

వృక్షంలోని ప్రతి \tetoken{చిహ్నిత సూత్రం}{!!{signed formula}}, వృక్షపు అగ్రభాగంలో జాబితా
చేసే ఒక \emph{పరికల్పన} అయినా అవుతుంది, లేదా దాని పైనున్న ఒక
\tetoken{చిహ్నిత సూత్రం}{!!a{signed formula}} నుంచి అనేక నిగమన నియమాల్లో ఒకదాని ద్వారా
పొందబడుతుంది. ముందున్న \tetoken{సూత్రానికి}{!!{formula}} ఉండగల ప్రతి \tetoken{ప్రధాన సంయోజకానికి}{!!{main operator}}
రెండు నియమాలు ఉంటాయి: సంకేతం~$\True$ అయిన సందర్భానికి ఒకటి,
సంకేతం~$\False$ అయిన సందర్భానికి ఒకటి. కొన్ని నియమాలు వృక్షాన్ని
శాఖలుగా విడగొడతాయి; మరికొన్ని శాఖకు \tetoken{చిహ్నిత సూత్రాలను}{!!{signed formula}s} మాత్రమే
చేర్చుతాయి. ఒక నియమాన్ని దానికి వెంటనే ముందున్న \tetoken{చిహ్నిత సూత్రానికి}{!!{signed formula}}
కాకుండా, మూలం నుంచి ఆ నియమం వర్తించే స్థానం వరకు శాఖలో ఉన్న ఏ
\tetoken{చిహ్నిత సూత్రానికైనా}{!!{signed formula}} వర్తింపజేయవచ్చు; తరచుగా అలా చేయాల్సిందే.

ఒక శాఖలో $\sFmla{\True}{!A}$, $\sFmla{\False}{!A}$ రెండూ ఉంటే అది
\emph{సంవృతమైనది}. ప్రతి శాఖా సంవృతంగా ఉంటే ఆ \tetoken{టాబ్లో}{!!{tableau}} సంవృతమైనది.
సహజార్థ వివరణలో ప్రతి శాఖ ఒక ఉమ్మడి సాధ్యతను వర్ణిస్తుంది; కానీ
$\sFmla{\True}{!A}$, $\sFmla{\False}{!A}$ రెండూ కలసి సాధ్యం కావు.
మరో మాటలో చెప్పాలంటే, ఒక శాఖ సంవృతమైతే అది వర్ణించే సాధ్యత
తొలగిపోయింది. ముఖ్యంగా, సంవృత \tetoken{టాబ్లో}{!!{tableau}} అనేది
$\sFmla{\True}{!A}$ రూపంలోని ప్రతి పరికల్పనను సత్యంగా, అలాగే
$\sFmla{\False}{!A}$ రూపంలోని ప్రతి పరికల్పనను అసత్యంగా ఒకేసారి
చేసే అన్ని సాధ్యతలనూ తొలగిస్తుంది.

$!A$ \emph{కోసం సంవృత \tetoken{టాబ్లో}{!!{tableau}}} అంటే మూలం~$\sFmla{\False}{!A}$గా
ఉండే సంవృత \tetoken{టాబ్లో}{!!{tableau}}. అటువంటి సంవృత \tetoken{టాబ్లో}{!!{tableau}} ఉంటే $!A$
అసత్యంగా ఉండే అన్ని సాధ్యతలూ తొలగిపోయాయి; అంటే ప్రతి
\tetoken{నిర్మాణంలోనూ}{!!{structure}} $!A$ సత్యం కావాలి.

\end{document}
